\documentclass[10pt,a4paper]{amsart}

\pdfinfoomitdate=1
\pdftrailerid{}
\pdfsuppressptexinfo=15

\usepackage[T1]{fontenc}
\usepackage{lmodern}
\usepackage[margin=24mm]{geometry}
\usepackage[protrusion=true,expansion=false]{microtype}
\usepackage{amsmath,amssymb,mathtools}
\usepackage{booktabs}
\usepackage{enumitem}
\usepackage{xcolor}
\usepackage[numbers,sort&compress]{natbib}
\usepackage[colorlinks=true,linkcolor=blue!55!black,citecolor=blue!55!black,urlcolor=blue!55!black]{hyperref}
\usepackage[nameinlink,noabbrev]{cleveref}

\setlist{nosep,leftmargin=*}
\raggedbottom
\allowdisplaybreaks
\emergencystretch=2em

\newtheorem{theorem}{Theorem}[section]
\newtheorem{proposition}[theorem]{Proposition}
\newtheorem{lemma}[theorem]{Lemma}
\newtheorem{corollary}[theorem]{Corollary}
\theoremstyle{definition}
\newtheorem{definition}[theorem]{Definition}
\theoremstyle{remark}
\newtheorem{remark}[theorem]{Remark}

\newcommand{\F}{\mathbb F}
\newcommand{\N}{\mathbb N}
\newcommand{\Mon}{\mathcal M}
\newcommand{\depth}{\operatorname{depth}}
\newcommand{\Fix}{\operatorname{Fix}}
\newcommand{\ord}{\operatorname{ord}}
\newcommand{\Res}{\mathcal U}

\title[Translation--GCD depth and fibres]{All-Depth Enumeration and Terminal Fibres\\
for Translation--GCD Erosion over Finite Fields}
\author{Anonymous}
\date{}
\subjclass[2020]{11T06, 05A15, 37P25, 12Y05}
\keywords{finite field, polynomial dynamics, translated gcd, cyclic run,
formal orbit Euler product, terminal fibre}

\begin{document}

\begin{abstract}
Let $q=p^a$ and iterate $T(f)=\gcd(f(x),f(x+1))$ on monic
polynomials over $\F_q$.  Translation has order $p$, so the process reaches
the invariant core $Q(f)=T^{p-1}(f)$.  We determine, for every
$0\leq t\leq p-1$, an exact all-degree generating function for the states
whose stabilization depth is at most $t$.  On each nonfixed orbit of
irreducible factors, subtracting the minimum exponent leaves a cyclic
nonnegative vector with a zero; its depth is its longest positive run.  A
finite run automaton supplies the local series, and a formal orbit Euler product over
irreducible translation orbits gives the global census.  We also prove a
unique graded split into an invariant core and the unit fibre $Q^{-1}(1)$.
Its generating function is $(1-qz^p)/(1-qz)$, yielding the exact-degree and
degree-capped fibre over every invariant target.  The translation-fixed
irreducible counts used in the formal orbit Euler product are due to Garefalakis and Reis, and
the underlying finite orbit-fold mechanism is treated as prior background;
the residual claims are the all-depth formal orbit Euler product and target-refined fibres.
Finite checks over $\F_4,\F_8$, and $\F_9$ are falsification controls only.
External release remains on hold.
\end{abstract}

\maketitle

\section{Scope and setup}

Let $q=p^a$ be a prime power and let $\Mon_q$ be the set of monic
polynomials in $\F_q[x]$, including $1$.  Write
\[
  \sigma f(x)=f(x+1),\qquad
  T(f)=\gcd(f,\sigma f),\qquad Q=T^{p-1},
\]
where every gcd is monic.  Since $\sigma^p=1$, the acting subgroup is the
prime-subfield copy of $\F_p$, not the full additive group of $\F_q$ when
$a>1$.  We use the stabilization depth
\[
 \depth(f)=\min\{t\geq0:T^t(f)=Q(f)\}.
\]
Thus fixed states have depth zero.

The invariant ring for this translation is
$\F_q[x^p-x]$.  Its monic elements have series
\begin{equation}\label{eq:invariant-series}
 I_{q,p}(z)=\sum_{h\in\Mon_q:\,\sigma h=h}z^{\deg h}
           =\frac{1}{1-qz^p},
\end{equation}
because $h=g(x^p-x)$ uniquely, with $g$ monic.  This characterization and
the translation-fixed irreducible classification are established in the
finite-field invariant literature, especially
\citet{Garefalakis2011} and \citet{Reis2018}; related group-action and
invariant-function frameworks include
\cite{Reimers2018,GowMcGuire2022,Schulz2023}.

We state the precise enumerative input.  Let
\begin{equation}\label{eq:Nd}
 N_d(q)=\frac1d\sum_{e\mid d}\mu(e)q^{d/e}
\end{equation}
be the number of degree-$d$ monic irreducibles.  Let $b_d$ count those
fixed by $\sigma$.  Then $b_d=0$ unless $d=pm$, and, if
$m=p^v s$ with $(s,p)=1$,
\begin{equation}\label{eq:bd}
 b_{pm}=\frac{p-1}{pm}\sum_{e\mid s}\mu(s/e)q^{p^v e}.
\end{equation}
This is exactly Reis's formula
\cite[Theorem~2(c)]{Reis2018}, after changing divisor variables, and is
\emph{not} a contribution here.  Consequently
\begin{equation}\label{eq:ad}
                  a_d=\frac{N_d(q)-b_d}{p}
\end{equation}
is the number of length-$p$ translation orbits of degree-$d$
irreducibles.  Equations \eqref{eq:Nd}--\eqref{eq:ad} are owned input to our
formal orbit Euler product.

The elementary identity
\begin{equation}\label{eq:window}
 T^t(f)=\gcd_{0\leq j\leq t}\sigma^j f
\end{equation}
follows by induction.  It implies $T^{p-1}(f)$ is invariant and
$\depth(f)\leq p-1$.  These facts, the sharp clock, the fixed set, and old
finite depth tables receive zero contribution credit.

\begin{remark}[P110 order-dual firewall]\label{rem:p110}
An earlier internal note, P110, studies $J(x)=x\vee\sigma x$ on a cyclically
acted lattice and obtains the consecutive orbit join, invariant endpoint,
finite clock, and recurrent points.  Equation \eqref{eq:window} is its
order-dual meet mechanism in the divisibility lattice.  We claim none of
those generic semilattice conclusions.  The contribution below begins only
with the polynomial-specific exponent census: cyclic run avoidance over
irreducible orbits and the resulting all-depth and target-fibre formulas.
\end{remark}

Shifted gcds and shift classes also occur in symbolic summation and
shiftless factorization \cite{GerhardEtAl2003}; those algorithmic
interfaces, like standard transfer matrices, are background rather than
novelty claims.

\section{Translation orbits and local depth}

Translation acts on monic irreducibles in orbits of length $1$ or $p$.
Fix a nonfixed orbit and choose an indexing
\[
 P_0,P_1=\sigma P_0,\ldots,P_{p-1}=\sigma^{p-1}P_0,
\]
with indices read modulo $p$.  If $e_i$ is the exponent of $P_i$ in $f$,
then, up to reversing the cyclic indexing, \eqref{eq:window} says that its
exponent after $t$ steps is
\begin{equation}\label{eq:sliding-min}
              e_i^{(t)}=\min_{0\leq j\leq t}e_{i+j}.
\end{equation}
At $t=p-1$ every coordinate equals $m=\min_i e_i$.  Put
$c_i=e_i-m$.  Then $c_i\geq0$ and $\min_i c_i=0$.

For a cyclic vector $c=(c_0,\ldots,c_{p-1})$ with a zero, let
$\lambda(c)$ denote the longest cyclic run of positive coordinates; set
$\lambda(0,\ldots,0)=0$.

\begin{lemma}[Local run statistic]\label{lem:run}
On a nonfixed irreducible orbit, the least $t$ for which the exponents in
\eqref{eq:sliding-min} have all reached their terminal value $m$ is
$\lambda(c)$.
\end{lemma}

\begin{proof}
After subtracting $m$, equation \eqref{eq:sliding-min} vanishes in every
coordinate exactly when every cyclic interval of length $t+1$ contains a
zero.  This is equivalent to the absence of a positive cyclic run of length
$t+1$, or to $\lambda(c)\leq t$.  The least such $t$ is
$\lambda(c)$, including the all-zero boundary.
\end{proof}

Fixed irreducibles contribute only to the invariant core and have no local
transient.  Hence the global depth of $f$ is the maximum of
$\lambda(c)$ over its nonfixed irreducible orbits, with the maximum of an
empty family equal to zero.

\section{The all-depth formal orbit Euler product}

For $0\leq t\leq p-1$, define the local residual series
\begin{equation}\label{eq:R-definition}
 R_{p,t}(y)=
 \sum_{\substack{c\in\N^p,\ \min_i c_i=0\\\lambda(c)\leq t}}
 y^{c_0+\cdots+c_{p-1}}.
\end{equation}
It has a finite transfer description.  Put $u=y/(1-y)$ and index rows and
columns by $0,\ldots,t$.  Let $M_t(u)$ be the $(t+1)\times(t+1)$ matrix
whose only nonzero entries are
\begin{equation}\label{eq:M}
 (M_t)_{i,0}=1\quad(0\leq i\leq t),\qquad
 (M_t)_{i,i+1}=u\quad(0\leq i<t).
\end{equation}

\begin{lemma}[Cyclic transfer]\label{lem:transfer}
For $0\leq t\leq p-1$,
\begin{equation}\label{eq:R-trace}
                         R_{p,t}(y)=\operatorname{tr}\bigl(M_t(u)^p\bigr).
\end{equation}
\end{lemma}

\begin{proof}
Read a cyclic support word of length $p$.  State $i$ records the current
positive run length.  A zero sends every state to $0$ with weight $1$; a
positive coordinate sends $i$ to $i+1$ with weight
$u=y+y^2+\cdots$, provided $i<t$.  A closed length-$p$ path is therefore a
cyclic word with no positive run longer than $t$, together with an arbitrary
positive height at every positive position.

Fix the cut at labelled coordinates $0,\ldots,p-1$.  The state before $j$
is uniquely forced to be the preceding positive-run length.  Starting after
any zero merely computes that same cyclic assignment, so multiple zeros do
not create paths.  Thus each labelled vector gives one closed path: the
trace neither quotients by rotations nor weights by zeros (and the all-zero
vector has its unique all-zero path).  Conversely every closed path recovers
this labelled support and its heights.  Summing closed paths is the trace in
\eqref{eq:R-trace}, and their weights are exactly those in
\eqref{eq:R-definition}.
\end{proof}

Write $[z^n]F(z)$ for coefficient extraction in formal power series.

\begin{theorem}[All-depth formal orbit Euler product]\label{thm:euler}
For every prime power $q=p^a$ and $0\leq t\leq p-1$, the ordinary
generating function for monic polynomials of depth at most $t$ is
\begin{equation}\label{eq:H}
 H_{q,p,t}(z)
 :=\sum_{f\in\Mon_q:\,\depth(f)\leq t}z^{\deg f}
 =\frac{1}{1-qz^p}
   \prod_{d\geq1}R_{p,t}(z^d)^{a_d},
\end{equation}
where $a_d$ is the owned input in \eqref{eq:ad} and $R_{p,t}$ is given by
\eqref{eq:R-trace}.  This formal orbit Euler product is coefficientwise well
defined.
\end{theorem}

\begin{proof}
Unique factorization partitions the irreducible factors of $f$ into fixed
orbits and length-$p$ orbits.  On each length-$p$ orbit, the exponent vector
splits uniquely as
\[
       (e_0,\ldots,e_{p-1})=m(1,\ldots,1)+c,
       \qquad \min_i c_i=0.
\]
All common-minimum pieces, together with every fixed irreducible factor,
form one arbitrary translation-invariant monic polynomial.  Their series is
$I_{q,p}(z)$ from \eqref{eq:invariant-series}.  For each of the $a_d$
nonfixed degree-$d$ orbits, \cref{lem:run} says that depth at most $t$ is
equivalent to $\lambda(c)\leq t$, and the residual degree is
$d\sum_i c_i$.  Its series is therefore $R_{p,t}(z^d)$.

The global depth is the maximum of the local depths, so these conditions
hold independently on every orbit.  Multiplication yields \eqref{eq:H}.
For a fixed coefficient $z^n$, only irreducible degrees $d\leq n$ can
contribute, proving formal well-definedness.
\end{proof}

\begin{corollary}[Exact depth layers]\label{cor:layers}
Set $H_{q,p,-1}(z)=0$.  For $0\leq t\leq p-1$,
\[
 \#\{f\in\Mon_q:\deg f=n,\ \depth(f)=t\}
   =[z^n]\bigl(H_{q,p,t}(z)-H_{q,p,t-1}(z)\bigr).
\]
Moreover,
\begin{equation}\label{eq:boundaries}
 R_{p,0}(y)=1,\quad H_{q,p,0}(z)=\frac1{1-qz^p},\qquad
 R_{p,p-1}(y)=\frac{1-y^p}{(1-y)^p},\quad
 H_{q,p,p-1}(z)=\frac1{1-qz}.
\end{equation}
\end{corollary}

\begin{proof}
The difference isolates the two nested depth thresholds.  At $t=0$, the
only normalized vector is zero.  At $t=p-1$, every nonnegative vector with
minimum zero is allowed, so its series is the series of all vectors minus
that of all-positive vectors:
\[
 (1-y)^{-p}-\bigl(y/(1-y)\bigr)^p
   =(1-y^p)/(1-y)^p.
\]
The first global identity in \eqref{eq:boundaries} is
\eqref{eq:invariant-series}; the last follows either from
\cref{thm:euler} and unique factorization or simply because every monic
polynomial has depth at most $p-1$ and there are $q^n$ of degree $n$.
\end{proof}

\section{Unit fibre and every invariant target}

The endpoint $Q(f)$ divides $f$ and is invariant.  Define the
\emph{unit fibre}
\[
                    \Res_{q,p}=Q^{-1}(1)
\]
and let $U_{q,p,n}$ count its degree-$n$ elements.  The word ``unit'' refers
only to the target $1$.

\begin{proposition}[Graded terminal split]\label{prop:split}
Multiplication gives a degree-preserving bijection
\begin{equation}\label{eq:split}
 \{h\in\Mon_q:\sigma h=h\}\times\Res_{q,p}
       \longrightarrow \Mon_q,\qquad (h,r)\longmapsto hr.
\end{equation}
Its inverse is $f\mapsto(Q(f),f/Q(f))$.  This is a set-theoretic graded
bijection, not a monoid quotient.
\end{proposition}

\begin{proof}
On every nonfixed irreducible orbit, $Q(f)$ has the common minimum exponent;
on a fixed irreducible it has the full exponent.  Hence $h=Q(f)$ is
invariant, divides $f$, and the residual $r=f/h$ has minimum exponent zero
on every nonfixed orbit and no fixed irreducible factor.  Thus $Q(r)=1$.
This proves existence and also forces both factors, proving uniqueness.

Equivalently, for invariant $h$ equation \eqref{eq:window} gives the
restricted identity
\begin{equation}\label{eq:restricted}
                    Q(hr)=hQ(r).
\end{equation}
This identity proves that every pair on the left of \eqref{eq:split} maps
back to itself.  It does not imply multiplicativity on arbitrary pairs.
\end{proof}

\begin{theorem}[Unit fibre and target fibres]\label{thm:fibres}
The unit-fibre generating function and coefficients are
\begin{align}
 U_{q,p}(z)&:=\sum_{n\geq0}U_{q,p,n}z^n
       =\frac{1-qz^p}{1-qz},\label{eq:U}\\
 U_{q,p,n}&=
 \begin{cases}
 q^n,&0\leq n<p,\\
 q^n-q^{n-p+1},&n\geq p.
 \end{cases}\label{eq:Un}
\end{align}
Let $h$ be any invariant monic polynomial of degree $m$.  Then, for every
$N\geq0$,
\begin{equation}\label{eq:exact-fibre}
 \#\{f\in\Mon_q:\deg f=N,\ Q(f)=h\}
 =\begin{cases}U_{q,p,N-m},&N\geq m,\\0,&N<m.\end{cases}
\end{equation}
Writing $L=D-m$, the degree-capped fibre through $D$ is zero for $L<0$;
for $L\geq0$ it equals
\begin{equation}\label{eq:capped}
 \sum_{n=0}^{L}U_{q,p,n}=
 \begin{cases}
 \dfrac{q^{L+1}-1}{q-1},&0\leq L<p,\\[6pt]
 \dfrac{q^{L+1}-1-q^{L-p+2}+q}{q-1},&L\geq p.
 \end{cases}
\end{equation}
Thus fibres depend on the target only through its degree.
\end{theorem}

\begin{proof}
All monic polynomials have series $(1-qz)^{-1}$.  Taking generating
functions in the graded bijection \eqref{eq:split} and using
\eqref{eq:invariant-series} gives
\[
       \frac1{1-qz}=\frac1{1-qz^p}U_{q,p}(z),
\]
which is \eqref{eq:U}; coefficient extraction gives \eqref{eq:Un}.

For a fixed invariant $h$, multiplication $r\mapsto hr$ maps the
degree-$(N-m)$ part of $\Res_{q,p}$ into the fibre over $h$ by
\eqref{eq:restricted}.  Conversely, \cref{prop:split} forces every element
of that fibre to be $hr$ for a unique such $r$.  This proves
\eqref{eq:exact-fibre}.  Summing \eqref{eq:Un} through $L$ gives
\eqref{eq:capped}; the second case subtracts
$q\sum_{j=0}^{L-p}q^j$ from $\sum_{j=0}^{L}q^j$.
\end{proof}

\begin{remark}[Why this is not a kernel]\label{rem:not-kernel}
The projection $Q$ is not a monoid homomorphism.  In characteristic $p$,
put $a=x$ and $b=(x^p-x)/x$.  The two exponent supports each miss a member
of the full translation orbit, so $Q(a)=Q(b)=1$, whereas
$ab=x^p-x$ is invariant and $Q(ab)=x^p-x$.  Hence $Q^{-1}(1)$ is a unit
fibre, not an algebraic kernel.  The proofs above use only orbit exponents
and the restricted identity \eqref{eq:restricted}.
\end{remark}

\section{Mechanical falsification and limitations}
\enlargethispage{3\baselineskip}

An independent verifier constructs
\[
 \F_4=\F_2[u]/(u^2+u+1),\quad
 \F_8=\F_2[u]/(u^3+u+1),\quad
 \F_9=\F_3[u]/(u^2+1)
\]
from explicit field tables.  It implements literal Horner translation,
Euclidean gcd, trial-division irreducibility, translation-orbit
partitioning, direct residual-vector enumeration, target fibres, and formal
orbit Euler product coefficients.  Separately, for $p=2,3$ and every $t$, it
constructs the truncated polynomial matrix $M_t(y/(1-y))$, takes the trace
of its $p$th power, and matches weights through $9$ against direct vector
enumeration.  It checks degrees $6,4,4$ in the three fields: $17{,}523$
states and $180{,}453$ assertions.  In particular, the $\F_9$ lane exercises the
intermediate depths $0,1,2$.  This finite enumeration is evidence against
small counterexamples, never a proof of
\cref{thm:euler,thm:fibres}.

The result is limited to the single order-$p$ translation generated by
$1$.  Characteristic zero has no such finite clock; simultaneous erosion by
a larger additive subspace is a different map.  We make no priority claim
for the literal dynamics, \eqref{eq:window}, the invariant ring, the fixed
irreducible formula, or the transfer method.  The owner-subtracted residual
is precisely the conjunction of \cref{thm:euler} with the target-refined
\cref{thm:fibres}.  A bounded literature non-hit is not novelty clearance,
and external posting or submission remains on hold pending specialist owner
review and explicit authorization.

\section{Conclusion}

After the known orbit fold and fixed-irreducible theory are removed, the
transient census is controlled by one local statistic: the longest positive
run in a normalized cyclic exponent vector.  Its transfer series produces
all depth thresholds at once, while the same normalization gives a uniform
graded fibre above every invariant endpoint.  These two formulas constitute
the complete internal claim package; the older clock, fixed set, and owner
formulas remain background.

\renewcommand{\bibfont}{\scriptsize}
\bibliographystyle{plainnat}
\bibliography{references}

\end{document}
